\section{Marco teórico}

En esta sección se introducen los conceptos mínimos necesarios para entender el resto del documento. El objetivo es tener un vocabulario común, no hacer desarrollos matemáticos largos.

\subsection{OpenMP: qué es y qué aporta}

OpenMP es una \textbf{API (Application Programming Interface) diseñada para la computación paralela en arquitecturas de memoria compartida}. No es un lenguaje ni una biblioteca: es un conjunto de \textbf{directivas de compilador}, \textbf{rutinas de biblioteca} y \textbf{variables de ambiente} que permiten al programador describir paralelismo sin escribir explícitamente la creación y sincronización de hilos.

El elemento central es la \textbf{región paralela}. Cuando el hilo maestro encuentra una región paralela, crea un \textit{equipo de hilos} (\textit{team}), todos ejecutan el bloque de código de la región, y al final de la región el hilo maestro espera a que todos terminen. Esa sincronización final es lo que garantiza que los resultados estén listos cuando la ejecución continúa.

La analogía con la versión de pthreads del caso de estudio 1 es directa: allí se escribía a mano el \texttt{pthread\_create} para cada hilo y el \texttt{pthread\_join} para esperarlos. OpenMP hace exactamente eso internamente, pero lo hace el \textit{runtime}.

Tres conceptos gobiernan el comportamiento de una región paralela:

\begin{itemize}
    \item \textbf{Directivas y centinelas}. Una directiva es una línea de código con significado \emph{solo} para el compilador. Se escriben con el centinela \texttt{\#pragma omp}. Un punto fundamental: \textbf{si el programa se compila sin la bandera \texttt{-fopenmp}, las directivas se ignoran silenciosamente} y el código se ejecuta secuencialmente. Por eso el flag de compilación es parte inseparable del uso de OpenMP.
    \item \textbf{Datos compartidos y privados}. En una región paralela, por defecto las variables son \emph{compartidas} (todos los hilos ven la misma copia) salvo las variables de control de los bucles, que son \emph{privadas} por defecto. El programador puede declarar explícitamente qué es compartido y qué es privado con las cláusulas \texttt{shared()} y \texttt{private()}.
    \item \textbf{Sincronización}. La región paralela tiene una \emph{barrier} implícita al final. Además existen construcciones explícitas (\texttt{critical}, \texttt{atomic}, \texttt{barrier}) para cuando varios hilos deben tocar el mismo dato.
\end{itemize}

\subsection{Los tres niveles de un programa paralelo}

En el caso de estudio 1 se implementó la paralelización manual, es decir, el nivel más bajo: el programador decide exactamente qué hace cada hilo. OpenMP introduce un nivel intermedio en el que el programador \emph{describe} qué bucles paralelizar y el \textit{runtime} decide el reparto. Existen además niveles más altos (\textit{tasking}, \textit{worksharing}) que no se usaron en este trabajo.

\subsection{Métricas de rendimiento: MIPS, FLOPS y por qué no sirven aquí}

El \textbf{rendimiento computacional} se puede medir de dos formas distintas:

\begin{itemize}
    \item \textbf{Tiempo de respuesta}: cuánto tarda el sistema en resolver \emph{una} aplicación concreta. Es la métrica del usuario final. Es la que se usa en este trabajo.
    \item \textbf{Throughput (productividad)}: cuánto trabajo es capaz de completar el sistema en un intervalo de tiempo. Es una métrica del \emph{sistema}, más cercana a la de un servidor con many workloads.
\end{itemize}

Para comparar dos máquinas se define el rendimiento como $R = 1/T$ y la diferencia porcentual:

\[
\%\Delta = \frac{R_x - R_y}{R_y} \times 100
\]

Las unidades de medida históricas para cuantificar el rendimiento computacional son:

\begin{itemize}
    \item \textbf{MIPS} (millions of instructions per second):Instructions per second. Es la unidad clásica, pero \textbf{depende fuertemente del repertorio de instrucciones} de cada microarquitectura, por lo que es poco comparable entre arquitecturas distintas.
    \item \textbf{FLOPS} (floating point operations per second): operaciones de coma flotante por segundo. Es la unidad moderna y se usa en el ranking \textit{Top500}. El problema es que \textbf{solo aplica a operaciones de coma flotante}.
\end{itemize}

En este trabajo el kernel trabaja con \textbf{enteros}, por lo que \textbf{no se puede reportar la aplicación en FLOPS}. Esta es una limitación real y no un descuido: la alternativa sería reescribir el kernel con \texttt{double}, lo que cambiaría el objeto de estudio (el algoritmo de punto flotante tiene sus propios problemas de precisión y de uso de la FPU) sin aportar nada a la pregunta que se quiere responder. Por eso se reportan \textbf{tiempos, speedup y diferencia porcentual de rendimiento}, que además es lo que el enunciado de la actividad pide de forma explícita.

\subsection{Benchmark: definición y tipos}

Un \textbf{benchmark} es un programa que permite evaluar el rendimiento de una arquitectura o de un equipo. Su lógica es generar una \emph{carga de trabajo} conocida y estandarizada para la cual se puede predecir el rendimiento en un trabajo real.

Según la procedencia de la carga de trabajo, los benchmarks se clasifican en:

\begin{itemize}
    \item \textbf{Programas reales completos}: por ejemplo un compilador, un procesador de textos, un outil CAD. Son los más representativos.
    \item \textbf{Kernels extraídos de programas reales}: es la categoría de este trabajo. Se extrae la parte central (el \textit{kernel}) de un programa real que es la que consume la mayor parte del tiempo de ejecución. El profesor insistió en medir el tiempo del \emph{kernel} de la multiplicación de matrices, no el programa completo, porque es el kernel lo que domina el tiempo y lo que determina el comportamiento de la caché.
    \item \textbf{Programas reducidos o \textit{toys}}: programas muy pequeños, de 10 a 100 líneas, que generan carga sin resolver un problema real.
    \item \textbf{Programas sintéticos}: creados artificialmente para generar carga en un recurso específico.
\end{itemize}

Una advertencia importante del profesor: los benchmarks reducidos y los sintéticos \textbf{no cargan la memoria principal de forma realista}, por lo que no son una buena opción para caracterizar el desempeño de una aplicación real. Para eso hay que usar programas reales o \emph{kernels} extraídos de problemas reales, como el nuestro.

\subsection{Perfilado: definición y tipos}

El \textbf{perfilado} (\textit{profiling}) tiene un objetivo distinto al benchmark. El benchmark mide el \emph{desempeño}; el perfilado busca \textbf{identificar el comportamiento del programa en una arquitectura específica}: dónde están los puntos ineficientes, dónde hay embotellamiento de memoria o de acceso a disco, y qué explica el escalamiento de la solución.

Las herramientas de perfilado se clasifican por su mecanismo:

\begin{itemize}
    \item \textbf{Contadores de hardware}: leen directamente los contadores del procesador (ciclos, instrucciones, fallos de caché).
    \item \textbf{Muestreo (\textit{sampling})}: interrumpen el programa en intervalos periódicos y capturan el estado, para después hacer un análisis estadístico. Es el mecanismo de \texttt{perf} y de \texttt{gprof}.
    \item \textbf{Trazado (\textit{tracing})}: en lugar de interrumpir periódicamente, detectan la ocurrencia de un evento concreto (por ejemplo, una llamada a función) y lo registran.
\end{itemize}

En este trabajo se usan dos herramientas, que cubren los dos ejes pedidos por la consigna:

\begin{itemize}
    \item \textbf{GProf}: perfilador por muestreo que se integra con el compilador. Se activa con la bandera \texttt{-pg}, que inserta instrumentación en el programa. Al ejecutarse, el programa escribe un archivo \texttt{gmon.out} con los contadores recolectados; luego la utilidad \texttt{gprof} convierte ese archivo binario en un reporte legible. Es la herramienta que el profesor presentó en clase.
    \item \textbf{Cachegrind}: perfilador de simulación de caché, parte de Valgrind. Reproduce el comportamiento de una jerarquía de caché (L1, LL) durante la ejecución y reporta, para cada función, cuántas instrucciones se ejecutaron y cuántas fallaron en caché. Con la opción \texttt{--cache-sim=yes} reporta además las tasas de fallo (\textit{miss rate}) de lectura y escritura.
\end{itemize}

\subsection{La jerarquía de caché y el miss rate}

El procesador no accede directamente a la memoria RAM: hay una jerarquía de cachés, pequeños y muy rápidos, entre el núcleo y la RAM.

\begin{center}
\begin{tabular}{|l|c|c|}
\hline
\textbf{Nivel} & \textbf{Tamaño típico} & \textbf{Tiempo de acceso} \\
\hline
L1 (datos e instrucciones) & 32--64 KiB por núcleo & ~4 ciclos \\
L2 & 512 KiB--2 MiB por núcleo & ~12 ciclos \\
L3 (compartida) & 8--32 MiB & ~40 ciclos \\
RAM & GiB & ~200 ciclos \\
\hline
\end{tabular}
\end{center}

El \textbf{miss rate} es la fracción de accesos a un nivel de caché que no se encuentran en ese nivel y deben bajar al nivel inferior. Es la métrica clave para entender el rendimiento de un kernel: un miss a RAM cuesta cientos de ciclos, mientras que un acierto a L1 cuesta unos pocos.

Cuando se dice que un acceso es \emph{contiguo} o \emph{secuencial}, se está describiendo el patrón con que se recorre un arreglo en memoria. En el algoritmo clásico de multiplicación de matrices, el acceso a $B$ no es contiguo, y esa es la raíz del problema que se mide en la sección 8.
